\section{2017级工科数学分析下 B卷}

\subsection{试题}

\paragraph*{一、填空题：共 5 题，每题 2 分，共 10 分。}
\begin{enumerate}
  \item 微分方程 $y''+2y'-3y=2-3x$ 的通解为\underline{\hspace{4cm}}。
  \item 向量场 $2xy\vec{i}+e^z\sin y\vec{j}+(x^2+y^2+z^2)\vec{k}$ 的旋度为\underline{\hspace{4cm}}。
  \item 设 $\Gamma$ 是平面上的下半圆周
  \[
    y=-\sqrt{1-x^2},\qquad -1\le x\le 1,
  \]
  则第一类曲线积分
  \[
    \int_\Gamma (x^2+y^2)\,ds=\underline{\hspace{4cm}}.
  \]
  \item 幂级数 $\displaystyle \sum_{n=1}^{\infty}(-1)^n\dfrac{1}{n\cdot 2^n}x^n$ 的收敛域为\underline{\hspace{4cm}}。
  \item 设周期为 $2\pi$ 的函数 $f(x)=|x|,\ -\pi<x\le\pi$，则 $f(x)$ 的傅里叶（Fourier）级数在 $x=\pi$ 处收敛于\underline{\hspace{4cm}}。
\end{enumerate}

\paragraph*{二、选择题：共 5 题，每题 2 分，共 10 分。}
\begin{enumerate}
  \item 下列微分方程中，属于二阶线性常微分方程的是（\quad）
  \begin{enumerate}
    \item[A.] $(x+y^2)\,dy+e^x\,dx=0$；
    \item[B.] $xy''+(\sin x)y'+(\ln x)y=\tan x$；
    \item[C.] $(y'')^2+y=2$；
    \item[D.] $y''+\ln y=2$。
  \end{enumerate}

  \item 二元函数 $f(x,y)$ 在点 $(a,b)$ 的两个偏导数
  \[
    \frac{\partial f}{\partial x}(a,b),\qquad
    \frac{\partial f}{\partial y}(a,b)
  \]
  存在是 $f(x,y)$ 在该点连续的（\quad）
  \begin{enumerate}
    \item[A.] 充分而非必要条件；
    \item[B.] 必要而非充分条件；
    \item[C.] 充分必要条件；
    \item[D.] 既非充分也非必要条件。
  \end{enumerate}

  \item 曲面 $3x^2+y^2+z^2=12$ 上点 $M(-1,0,3)$ 处的切平面与平面 $z=0$ 的夹角是（\quad）
  \begin{enumerate}
    \item[A.] $\dfrac{\pi}{6}$；
    \item[B.] $\dfrac{\pi}{4}$；
    \item[C.] $\dfrac{\pi}{3}$；
    \item[D.] $\dfrac{\pi}{2}$。
  \end{enumerate}

  \item 设 $\Gamma$ 是平面曲线
  \[
    \frac{x^2}{a^2}+\frac{y^2}{b^2}=1,\qquad a,b>0,
  \]
  取逆时针方向，则第二类曲线积分
  \[
    \int_\Gamma x\,dy-y\,dx
  \]
  等于（\quad）
  \begin{enumerate}
    \item[A.] $\pi ab$；
    \item[B.] $2\pi ab$；
    \item[C.] $0$；
    \item[D.] $-2\pi ab$。
  \end{enumerate}

  \item 下列级数发散的是（\quad）
  \begin{enumerate}
    \item[A.] $\displaystyle \sum_{n=1}^{\infty}2^n\sin\frac{\pi}{3^n}$；
    \item[B.] $\displaystyle \sum_{n=2}^{\infty}\frac{1}{\ln n}$；
    \item[C.] $\displaystyle \sum_{n=2018}^{\infty}\frac{2^n}{n!}$；
    \item[D.] $\displaystyle \sum_{n=1}^{\infty}\left(1-\frac1n\right)^{n^2}$。
  \end{enumerate}
\end{enumerate}

\paragraph*{三、计算题：共 3 题，每题 10 分，共 30 分。}
\begin{enumerate}
  \item 设
  \[
    z=xf\left(\frac{y}{x}\right)+g(x,xy),
  \]
  其中 $f(t)$ 具有连续的二阶导数，$g(u,v)$ 具有连续的二阶偏导数，计算二阶偏导数
  \[
    \frac{\partial^2 z}{\partial x\partial y}.
  \]

  \item 计算累次积分
  \[
    \int_{\frac14}^{\frac12}dx\int_{\frac12}^{\sqrt{x}}\cos\left(\pi\frac{x}{y}\right)\,dy
    +\int_{\frac12}^{1}dx\int_x^{\sqrt{x}}\cos\left(\pi\frac{x}{y}\right)\,dy.
  \]

  \item 计算曲面 $z=\sqrt{x^2+y^2}$ 夹在两柱面
  \[
    x^2+y^2=x,\qquad x^2+y^2=2x
  \]
  之间的那一部分的面积。
\end{enumerate}

\paragraph*{四、解答题：共 3 题，每题 10 分，共 30 分。}
\begin{enumerate}
  \item 设 $\Sigma$ 是曲线
  \[
    \begin{cases}
      z=\sqrt{y-1},\\
      x=0,
    \end{cases}
    \qquad 1\le y\le 3
  \]
  绕 $y$ 轴旋转一周所成的曲面，其法向量与 $y$ 轴正向夹角恒大于 $\dfrac{\pi}{2}$。求曲面积分
  \[
    \iint_\Sigma (8y+1)x\,dy\,dz+2(1-y^2)\,dz\,dx+(z^2-2z)\,dx\,dy.
  \]

  \item 设曲线积分
  \[
    \int_\Gamma (f(x)-e^x)\sin\frac{y}{2}\,dx-\frac12 f(x)\cos\frac{y}{2}\,dy
  \]
  与路径无关，其中 $f(x)$ 有一阶连续导数且 $f(0)=0$。求 $f(x)$，并计算曲线积分
  \[
    \int_{(0,0)}^{(1,1)}(f(x)-e^x)\sin\frac{y}{2}\,dx-\frac12 f(x)\cos\frac{y}{2}\,dy.
  \]

  \item 将函数 $f(x)=\cos^2x$ 展开成 $x$ 的幂级数。
\end{enumerate}

\paragraph*{五、证明题：共 1 题，每题 10 分，共 10 分。}
设 $0<c<1$，证明函数项级数
\[
  \sum_{n=0}^{\infty}(1-x)^n
\]
在 $[c,1]$ 一致收敛，但在 $(0,1)$ 不一致收敛。

\paragraph*{六、应用题：共 1 题，每题 10 分，共 10 分。}
设椭圆
\[
  x^2+3y^2=12
\]
的内接等腰三角形之底边平行于椭圆的长轴，求这样的三角形的最大面积。

\subsection{答案与解析}

\paragraph*{一、填空题答案}
\begin{enumerate}
  \item
  \[
    y=C_1e^{-3x}+C_2e^x+x,\qquad C_1,C_2\text{ 为任意常数}.
  \]

  \item 由旋度公式
  \[
    \nabla\times\vec{F}
    =
    \begin{vmatrix}
      \vec{i} & \vec{j} & \vec{k}\\
      \partial_x & \partial_y & \partial_z\\
      2xy & e^z\sin y & x^2+y^2+z^2
    \end{vmatrix}
    =(2y-e^z\sin y)\vec{i}-2x\vec{j}-2x\vec{k}.
  \]

  \item 因 $\Gamma$ 为单位圆下半圆周，故 $x^2+y^2=1$，弧长为 $\pi$，所以
  \[
    \int_\Gamma (x^2+y^2)\,ds=\int_\Gamma 1\,ds=\pi.
  \]

  \item 答案页给出的收敛域为
  \[
    (-2,2].
  \]
  若按有意义的幂级数下标 $n=1$ 起算，则半径为 $2$，在 $x=2$ 为交错调和级数收敛，在 $x=-2$ 为调和级数发散，故收敛域为 $(-2,2]$。


  \item 傅里叶级数在 $x=\pi$ 处收敛于左右极限平均值。由周期延拓得
  \[
    \frac{f(\pi-0)+f(\pi+0)}{2}=\frac{\pi+\pi}{2}=\pi.
  \]
\end{enumerate}

\paragraph*{二、选择题}
\begin{enumerate}
  \item 只有 \(xy''+(\sin x)y'+(\ln x)y=\tan x\) 对未知函数 $y$ 及其导数均为一次，是二阶线性常微分方程。选 B。
  \item 偏导数存在不能推出连续，连续也不能推出偏导数存在，因此二者既非充分也非必要。选 D。
  \item 曲面法向量为 \((6x,2y,2z)\)，在 $(-1,0,3)$ 处为 $(-6,0,6)$；它与平面 $z=0$ 的法向量 $(0,0,1)$ 夹角为 $\pi/4$。选 B。
  \item 由 Green 公式，\(\displaystyle \oint_\Gamma x\,dy-y\,dx=2S=2\pi ab\)。选 B。
  \item A、C、D 均收敛；\(\displaystyle \sum_{n=2}^{\infty}\frac1{\ln n}\) 发散。选 B。
\end{enumerate}

\paragraph*{三、计算题答案与解析}
\begin{enumerate}
  \item 先对 $y$ 求偏导：
  \[
    \frac{\partial z}{\partial y}
    =f'\left(\frac{y}{x}\right)+xg_2'(x,xy).
  \]
  再对 $x$ 求偏导，得
  \[
    \frac{\partial^2z}{\partial x\partial y}
    =
    -\frac{y}{x^2}f''\left(\frac{y}{x}\right)
    +g_2'(x,xy)+xg_{21}''(x,xy)+xyg_{22}''(x,xy).
  \]

  \item 记 $D$ 为由 $y=\dfrac12,\ y=\sqrt{x}$ 和 $y=x$ 围成的区域。原积分等于
  \begin{align*}
    &\int_{\frac14}^{\frac12}dx\int_{\frac12}^{\sqrt{x}}\cos\left(\pi\frac{x}{y}\right)\,dy
    +\int_{\frac12}^{1}dx\int_x^{\sqrt{x}}\cos\left(\pi\frac{x}{y}\right)\,dy\\
    &\qquad
    =\iint_D\cos\left(\pi\frac{x}{y}\right)\,dx\,dy
    =\int_{\frac12}^{1}dy\int_{y^2}^{y}\cos\left(\pi\frac{x}{y}\right)\,dx.
  \end{align*}
  因此
  \begin{align*}
    \int_{\frac12}^{1}dy\int_{y^2}^{y}\cos\left(\pi\frac{x}{y}\right)\,dx
    &=-\frac1\pi\int_{\frac12}^{1}y\sin(\pi y)\,dy\\
    &=\frac{1-\pi}{\pi^3}.
  \end{align*}

  \item 曲面的面积微元为
  \begin{align*}
    dS
    &=\sqrt{1+\left(\frac{\partial z}{\partial x}\right)^2+\left(\frac{\partial z}{\partial y}\right)^2}\,dx\,dy\\
    &=\sqrt{1+\left(\frac{x}{\sqrt{x^2+y^2}}\right)^2+\left(\frac{y}{\sqrt{x^2+y^2}}\right)^2}\,dx\,dy\\
    &=\sqrt2\,dx\,dy.
  \end{align*}
  投影区域为 $x\le x^2+y^2\le 2x$，其面积为
  \[
    \pi-\frac14\pi=\frac34\pi.
  \]
  故所求面积为
  \[
    A=\iint_{x\le x^2+y^2\le 2x}\sqrt2\,dx\,dy
    =\sqrt2\left(\pi-\frac14\pi\right)
    =\frac{3\sqrt2}{4}\pi.
  \]
\end{enumerate}

\paragraph*{四、解答题答案与解析}
\begin{enumerate}
  \item 设 $\Sigma_1$ 为平面 $y=3$ 在 $x^2+z^2\le 2$ 的部分，取法向量 $(0,1,0)$。则 $\Sigma$ 与 $\Sigma_1$ 形成闭曲面，外侧为正向，所围区域为 $\Omega$。由 Gauss 公式，
  \begin{align*}
    &\iint_{\Sigma\cup\Sigma_1}(8y+1)x\,dy\,dz+2(1-y^2)\,dz\,dx+(z^2-2z)\,dx\,dy\\
    &\qquad
    =\iiint_\Omega (4y+2z-1)\,dx\,dy\,dz.
  \end{align*}
  由关于 $z$ 的对称性，
  \[
    \iiint_\Omega z\,dx\,dy\,dz=0.
  \]
  因此
  \begin{align*}
    \iiint_\Omega (4y+2z-1)\,dx\,dy\,dz
    &=\iiint_\Omega (4y-1)\,dx\,dy\,dz\\
    &=\int_1^3dy\iint_{x^2+z^2\le y-1}(4y-1)\,dx\,dz\\
    &=\int_1^3(4y-1)\pi(y-1)\,dy
    =\frac{50}{3}\pi.
  \end{align*}
  在 $\Sigma_1$ 上，$Q=2(1-y^2)=-16$，故
  \[
    \iint_{\Sigma_1}(8y+1)x\,dy\,dz+2(1-y^2)\,dz\,dx+(z^2-2z)\,dx\,dy
    =\iint_{x^2+z^2\le2}(-16)\,dx\,dz=-32\pi.
  \]
  所以
  \[
    \iint_\Sigma (8y+1)x\,dy\,dz+2(1-y^2)\,dz\,dx+(z^2-2z)\,dx\,dy
    =\frac{50}{3}\pi+32\pi
    =\frac{146}{3}\pi.
  \]

  \item 记
  \[
    P=(f(x)-e^x)\sin\frac{y}{2},\qquad
    Q=-\frac12f(x)\cos\frac{y}{2}.
  \]
  由于曲线积分与路径无关，故
  \[
    \frac{\partial P}{\partial y}=\frac{\partial Q}{\partial x}.
  \]
  即
  \[
    f(x)-e^x=-f'(x).
  \]
  解一阶线性方程 $f'(x)+f(x)=e^x$，得
  \[
    f(x)=\frac12e^x+Ce^{-x}.
  \]
  由 $f(0)=0$，得 $C=-\dfrac12$，所以
  \[
    f(x)=\frac12e^x-\frac12e^{-x}.
  \]
  此时
  \begin{align*}
    P\,dx+Q\,dy
    &=-\frac12(e^x+e^{-x})\sin\frac{y}{2}\,dx
      -\frac14(e^x-e^{-x})\cos\frac{y}{2}\,dy\\
    &=d\left[-\frac12(e^x-e^{-x})\sin\frac{y}{2}\right].
  \end{align*}
  因此
  \begin{align*}
    \int_{(0,0)}^{(1,1)}P\,dx+Q\,dy
    &=
    \left[-\frac12(e^x-e^{-x})\sin\frac{y}{2}\right]_{(0,0)}^{(1,1)}\\
    &=-\frac12(e-e^{-1})\sin\frac12.
  \end{align*}

  \item 由
  \[
    \cos t=\sum_{n=0}^{\infty}\frac{(-1)^n}{(2n)!}t^{2n},
  \]
  以及
  \[
    \cos^2x=\frac12+\frac12\cos2x,
  \]
  得
  \begin{align*}
    \cos^2x
    &=\frac12+\frac12\sum_{n=0}^{\infty}\frac{(-1)^n}{(2n)!}(2x)^{2n}\\
    &=\frac12+\frac12\sum_{n=0}^{\infty}\frac{(-1)^n4^n}{(2n)!}x^{2n}.
  \end{align*}
  其收敛域为 $(-\infty,+\infty)$。
\end{enumerate}

\paragraph*{五、证明题答案与解析}
证明：对任意 $x\in[c,1]$，有
\[
  (1-x)^n\le (1-c)^n.
\]
由于级数
\[
  \sum_{n=0}^{\infty}(1-c)^n
\]
收敛，由 Weierstrass 优级数判别法，函数项级数
\[
  \sum_{n=0}^{\infty}(1-x)^n
\]
在 $[c,1]$ 上一致收敛。

另一方面，该级数在 $(0,1]$ 上逐点收敛。其前 $n+1$ 项部分和为
\[
  S_n(x)=\sum_{k=0}^{n}(1-x)^k
  =\frac{1-(1-x)^{n+1}}{x},
\]
和函数为
\[
  S(x)=\frac1x.
\]
取 $x_n=\dfrac1n\in(0,1)$，则
\[
  \left|S_n(x_n)-S(x_n)\right|
  =\left(1-\frac1n\right)^n(n-1)\to+\infty.
\]
因此函数项级数
\[
  \sum_{n=0}^{\infty}(1-x)^n
\]
在 $(0,1)$ 不一致收敛。

\paragraph*{六、应用题答案与解析}
设等腰三角形的三个顶点为 $(-x,y)$、$(x,y)$ 和 $(0,2)$，其中 $x>0$。则三角形面积为
\[
  f(x,y)=x(2-y).
\]
只需求函数 $f(x,y)$ 在条件 $x^2+3y^2=12$ 下的条件极值。设 Lagrange 函数
\[
  L(x,y,\lambda)=x(2-y)-\lambda(x^2+3y^2-12).
\]
极值点满足
\[
  \begin{cases}
    \dfrac{\partial L}{\partial x}=2-y-2\lambda x=0,\\
    \dfrac{\partial L}{\partial y}=-x-6\lambda y=0,\\
    \dfrac{\partial L}{\partial \lambda}=-(x^2+3y^2-12)=0.
  \end{cases}
\]
解得
\[
  (x,y,\lambda)=\left(3,-1,\frac12\right),\quad
  \left(-3,-1,-\frac12\right),\quad
  (0,2,0).
\]
因 $x>0$，极大值点为 $(x,y)=(3,-1)$。因此三角形最大面积为
\[
  f(3,-1)=9.
\]
